% =====================================================================
%  MIC Summer-of-Code  --  Week 5
%  Neural Networks for Images (CNNs, U-Net, and the modern zoo)
%
%  Self-contained handbook chapter.  Compile with pdflatex (twice) or
%  upload to Overleaf and hit Recompile.
%      pdflatex MIC_Neural_Networks_for_Images.tex
%      pdflatex MIC_Neural_Networks_for_Images.tex
%
%  Built from Prof. Suyash P. Awate's CS-736 DeepGenerativeModeling deck
%  (CNN half), with attribution, for the use of mentees on this project.
% =====================================================================
\documentclass[11pt]{article}

\usepackage[a4paper,margin=1in]{geometry}
\usepackage{amsmath,amssymb,amsthm,mathtools,bm}
\usepackage{enumitem}
\usepackage{booktabs}
\usepackage{xcolor}
\usepackage{microtype}
\usepackage[most]{tcolorbox}
\usepackage{titlesec}
\usepackage[colorlinks=true,linkcolor=blue!45!black,urlcolor=blue!55!black,
            citecolor=blue!45!black]{hyperref}

\titleformat{\section}{\Large\bfseries\sffamily\color{blue!40!black}}{\thesection}{0.6em}{}
\titleformat{\subsection}{\large\bfseries\sffamily\color{blue!32!black}}{\thesubsection}{0.5em}{}

\newtcolorbox{keybox}[1][]{colback=blue!4!white,colframe=blue!40!black,
  fonttitle=\bfseries\sffamily,title=#1,boxrule=0.6pt,arc=2pt,left=4pt,right=4pt,top=3pt,bottom=3pt}
\newtcolorbox{notebox}[1][Aside]{colback=black!3!white,colframe=black!45!white,
  fonttitle=\bfseries\sffamily,title=#1,boxrule=0.5pt,arc=2pt,left=4pt,right=4pt,top=3pt,bottom=3pt}

\newcommand{\R}{\mathbb{R}}
\newcommand{\norm}[1]{\left\lVert #1 \right\rVert}
\newcommand{\deck}[1]{\textsf{\small[#1]}}
\DeclareMathOperator*{\argmin}{arg\,min}

\title{\textbf{Neural Networks for Images}\\[2pt]
\large Convolutional networks, U-Net, and the modern segmentation zoo\\[6pt]
\normalsize Medical Image Computing -- Summer of Code -- Week 5}
\author{}
\date{}

\begin{document}
\maketitle
\vspace{-2.2em}

\begin{keybox}[How to read this chapter]
For four weeks you \emph{hand-built} the pieces of an imaging algorithm: a noise
model, a prior, a forward operator, a distance. Deep learning replaces the
hand-built pieces with \emph{learned} ones --- but it does not throw away the
Bayesian picture, it slots into it. A CNN is just a very flexible function whose
millions of parameters are fitted by gradient descent; in our language it is
``one class of priors and one class of optimisers''. This chapter builds the CNN
from the neuron up, explains why convolution is the right inductive bias for
images, tours the architectures you will actually meet in medical imaging
(U-Net, nnU-Net, ViT, SAM), and is honest about what deep learning does and does
not buy you. Keep the \deck{DeepGenerativeModeling} deck open (the CNN material
is slides 4--29); do the notebook to feel the learning mechanics in NumPy. The
heavy PyTorch practice lives in the cited tutorials --- this chapter is the map.
\end{keybox}

\tableofcontents
\bigskip

% =====================================================================
\section{From hand-built features to learned features}
% =====================================================================

Weeks 1--4 were a course in \emph{designing} the right function for a task: a
Huber penalty that preserves edges, a ramp filter that inverts the Radon
transform, a Procrustes alignment that removes pose. Each was a clever, fixed
recipe. Deep learning asks a different question: \emph{what if we let the data
choose the recipe?} A neural network is a function $f_\theta$ with a huge number
of tunable parameters $\theta$; we pick a loss that measures how wrong it is on
training examples, and we let gradient descent adjust $\theta$ until it is right.

This is not a betrayal of everything before --- it is the same machinery with the
function class swapped out \deck{DeepGenerativeModeling, slide 24}. You still
choose a likelihood (the loss), you still regularise (weight decay, dropout,
augmentation), and you still optimise (gradient descent, now with automatic
differentiation). The course's stance is deliberately un-hyped: deep learning is
\emph{one} powerful family of models and \emph{one} family of optimisers, to be
understood inside the probabilistic picture you already have --- not a magic box.

% =====================================================================
\section{The neuron, the layer, and why nonlinearity}
% =====================================================================

The atom is the \textbf{artificial neuron}: take inputs $x$, form a weighted sum
$w^{\!\top}x + b$, and pass it through a nonlinear \textbf{activation} $\sigma$:
\[
a = \sigma\!\left(w^{\!\top}x + b\right).
\]
Stack many neurons into a \textbf{layer} (a matrix multiply plus a bias plus a
pointwise nonlinearity), and stack many layers into a \textbf{network}
\deck{DeepGenerativeModeling, slides 4, 13}. The depth is where the power comes
from: each layer transforms the representation produced by the previous one, so
early layers can learn simple patterns and later layers combine them into
complex ones.

\begin{keybox}[Why the nonlinearity is not optional]
Without $\sigma$, a stack of layers is just $W_3(W_2(W_1 x))=(W_3W_2W_1)x$ --- one
big linear map, no matter how deep. The nonlinear activation between layers is
\emph{exactly} what lets a deep network represent functions a single linear layer
cannot. ``Deep'' only means something because of the nonlinearity in between.
\end{keybox}

% =====================================================================
\section{Convolution: the right inductive bias for images}
% =====================================================================

A fully-connected layer on a $256\times256$ image would need a weight for every
pixel-to-neuron pair --- billions of parameters, and it would have to re-learn
that an edge is an edge separately at every location. Images have structure we
can build in instead. A \textbf{convolution layer} slides a small learnable
\textbf{filter} (kernel) across the image and records, at each position, how
strongly the local patch matches the filter \deck{DeepGenerativeModeling, slides
5--7}. Three properties make this the right choice for images:
\begin{itemize}[nosep,leftmargin=*]
  \item \textbf{Locality:} each output looks only at a small neighbourhood ---
  the receptive field --- because image structure is local.
  \item \textbf{Weight sharing:} the same filter is used everywhere, so a feature
  (say, a vertical edge) is detected wherever it occurs, and the parameter count
  is tiny (the filter size, not the image size).
  \item \textbf{Translation equivariance:} shift the input, and the feature map
  shifts the same way --- the network does not need to see an object in every
  position to recognise it.
\end{itemize}
Stacked convolution layers grow the receptive field and build a hierarchy: the
first layer of a trained network learns oriented edges and colour blobs; the
second learns corners and textures; deeper layers learn object parts. This is not
a metaphor --- you can visualise the learned filters of AlexNet and literally see
the edge detectors emerge in layer 1 \deck{DeepGenerativeModeling, slides
9--10}. Strikingly, the same hierarchy shows up in the visual cortex (simple
cells $\to$ complex cells), which is where the idea came from
\deck{DeepGenerativeModeling, slides 16--21}.

\begin{notebox}[``Convolution'' in deep learning is really cross-correlation]
What DL libraries call convolution does \emph{not} flip the kernel --- it is
mathematical \textbf{cross-correlation} \deck{DeepGenerativeModeling, slide 8}.
For \emph{learning} it makes no difference (the network just learns the flipped
filter), but in the imaging literature ``convolution'' means convolution and
``cross-correlation'' means cross-correlation. Keep the two vocabularies
straight: it is the one place the same word means two things across the weeks of
this course.
\end{notebox}

\paragraph{The knobs.} A convolution layer is specified by its kernel size, the
number of input and output \textbf{channels}, the \textbf{stride} (how far the
filter hops --- stride $>1$ downsamples), the \textbf{padding} (zeros around the
border, to control output size), and the \textbf{dilation} (gaps in the kernel,
to enlarge the receptive field cheaply). The parameter count of one layer is
\[
(\text{kernel height}\times\text{kernel width}\times C_{\text{in}} + 1)\times C_{\text{out}},
\]
which for a $3\times3$ kernel mapping $64\to64$ channels is
$(3\cdot3\cdot64+1)\cdot64 \approx 37$k --- independent of the image size.

\begin{notebox}[Worked example: why convolution saves parameters]
Take a $256\times256$ grayscale image. A \emph{fully-connected} layer with $1000$
hidden units needs $256^2\times1000\approx 6.6\times10^7$ weights --- and it would
learn ``edge at top-left'' separately from ``edge at bottom-right''. A
\emph{convolution} layer with $64$ filters of size $3\times3$ needs only
$(3\cdot3+1)\cdot64=640$ weights and detects an edge \emph{anywhere}. Stacking
layers grows the \textbf{receptive field} (the input region one output sees):
two $3\times3$ layers see a $5\times5$ patch, three see $7\times7$. So depth, not
width, is how a CNN assembles global understanding from cheap local filters.
\end{notebox}

% =====================================================================
\section{The supporting cast: activations, pooling, normalisation}
% =====================================================================

A working network needs a few more standard layers \deck{DeepGenerativeModeling,
slides 11--13}:

\begin{description}[leftmargin=*,nosep]
  \item[\textbf{Activations.}] \textbf{ReLU} $f(x)=\max(0,x)$ is the default ---
  cheap, and it does not saturate for positive inputs (so gradients flow).
  \textbf{GELU} $f(x)=x\,\Phi(x)$ (a smooth ReLU) is common in modern nets;
  \textbf{softplus} $\log(1+e^x)$ is a smooth positive map; \textbf{sigmoid} and
  \textbf{softmax} squash outputs into $[0,1]$ / a probability vector for the
  \emph{final} classification layer.
  \item[\textbf{Pooling.}] \textbf{Max-} or \textbf{average-pooling} shrinks a
  feature map (e.g.\ $2\times2\to1$), giving a little translation invariance and
  cutting compute. Modern nets often replace pooling with strided convolutions.
  \item[\textbf{Normalisation.}] \textbf{BatchNorm} standardises activations
  across the batch to stabilise and speed training; for the \emph{small batches}
  typical of 3D medical data, \textbf{GroupNorm} or \textbf{InstanceNorm} are
  preferred because BatchNorm's statistics get noisy.
\end{description}

% =====================================================================
\section{Training: loss, backprop, optimisers, overfitting}
% =====================================================================

\paragraph{The loss.} The loss turns ``how wrong is the network'' into a number
to minimise. For \textbf{classification} it is usually \textbf{cross-entropy}
(the negative log-likelihood of a categorical model). For \textbf{segmentation}
--- where one class often dwarfs another (a tiny tumour in a big scan) --- plain
accuracy and cross-entropy are misleading, so people use the \textbf{Dice loss}
(or generalised Dice / Tversky), which directly rewards overlap with the target
mask. Choosing the loss \emph{is} choosing the likelihood: it is the same modelling
decision you made when you picked Gaussian vs Rician in Week 2.

\paragraph{Backpropagation.} How do we get the gradient of the loss with respect
to millions of parameters? \textbf{Reverse-mode automatic differentiation}:
the network is a graph of simple operations, each of which knows its local
derivative; backprop applies the chain rule from the loss backwards through the
graph, reusing intermediate results, so the \emph{whole} gradient costs about as
much as one forward pass \deck{DeepGenerativeModeling, slide 3}. You almost never
write it by hand --- PyTorch's autograd does it --- but you will implement a tiny
version in the notebook so it is not magic.

\begin{notebox}[Backprop in one concrete line]
For a two-layer net $a=\sigma(W_1x)$, $\hat y=W_2 a$, the chain rule gives
\[
\frac{\partial L}{\partial W_2}=\frac{\partial L}{\partial\hat y}\,a^{\!\top},
\qquad
\frac{\partial L}{\partial W_1}=\Big(W_2^{\!\top}\tfrac{\partial L}{\partial\hat y}\odot\sigma'(W_1x)\Big)x^{\!\top}.
\]
That is the \emph{whole} algorithm: push the error backwards, multiplying by each
layer's local derivative ($\sigma'$ at a nonlinearity, a weight matrix at a linear
layer). The notebook codes exactly this for a 1-hidden-layer net.
\end{notebox}

\paragraph{Optimisers.} Plain \textbf{SGD} updates $\theta\leftarrow\theta-\eta\nabla
L$; \textbf{momentum} smooths the trajectory; \textbf{Adam}/\textbf{AdamW} adapt
a per-parameter step size and are the usual default. A \textbf{learning-rate
schedule} (warm-up then decay) matters more than people expect.

\paragraph{Overfitting and how to fight it.} A network with millions of
parameters can simply \emph{memorise} a small dataset. The defences are the
regularisers you already know in a new guise \deck{DeepGenerativeModeling, slide
29}: \textbf{weight decay} (an $\ell^2$ prior on the weights --- literally
Tikhonov), \textbf{dropout} (randomly zero activations during training),
\textbf{data augmentation} (rotate/flip/deform training images so the model sees
more variety), \textbf{early stopping}, and getting more data. In medical imaging,
where 100 labelled scans is ``a lot'', these are not optional --- they are the
difference between a model that generalises and one that does not.

\subsection{Surviving small data: augmentation, transfer, self-supervision}

Three tools matter so much in medical imaging that they deserve their own
heading. \textbf{Data augmentation} manufactures new training examples with
label-preserving transforms --- flips, rotations, intensity jitter, and (very
medical) \emph{elastic deformations} that mimic anatomical variability; it is the
cheapest, most reliable regulariser when labels are scarce. \textbf{Transfer
learning} starts from a network pre-trained on a large dataset (ImageNet, or a
medical foundation model) and fine-tunes it on your few hundred scans --- the
early edge/texture filters transfer almost for free. \textbf{Self-supervised
pre-training} goes one better: learn features from \emph{unlabelled} scans (by
inpainting, contrastive learning, or masked-image modelling), then fine-tune on
the tiny labelled set. Together, these are why a 200-scan dataset can still train
a usable model --- and why ``just collect more labels'' is rarely the first move.

% =====================================================================
\section{Architectures you will actually meet}
% =====================================================================

\paragraph{Classifiers (the lineage).} \textbf{LeNet} (1998) and \textbf{AlexNet}
(2012) established the conv$\to$pool$\to$dense pattern; \textbf{VGG} made it deep
and uniform \deck{DeepGenerativeModeling, slide 14}. \textbf{ResNet} added
\textbf{residual (skip) connections} $y=f(x)+x$, which let gradients flow through
very deep nets and made 100+ layer networks trainable --- the single most
important architectural trick of the 2010s.

\paragraph{Encoder--decoder and U-Net.} For tasks whose \emph{output is an image}
(segmentation, denoising, super-resolution) you need an output the same size as
the input. An \textbf{encoder--decoder} compresses the image to a small latent
representation and then expands it back \deck{DeepGenerativeModeling, slides
23--27}. The decisive refinement for medical imaging is \textbf{U-Net} (MICCAI
2015): \textbf{skip connections} that concatenate encoder feature maps to the
decoder at each scale, so fine spatial detail lost during downsampling is handed
back during upsampling \deck{DeepGenerativeModeling, slide 26}. U-Net is the
\emph{de facto} backbone of medical image segmentation, with the original paper
passing 100k+ citations.

\paragraph{nnU-Net: configuration beats novelty.} \textbf{nnU-Net} (``no-new-Net'')
is not a new architecture --- it is a self-configuring \emph{pipeline} that, given
a dataset, automatically sets preprocessing, patch size, normalisation, and
training schedule around a fairly plain U-Net. Its lesson is humbling and
important: careful configuration of an old architecture beats most clever new
ones, and nnU-Net remains the benchmark to beat in 3D medical segmentation.

\paragraph{Transformers and foundation models.} \textbf{Vision Transformers}
(ViT) replace convolution's locality with global \emph{self-attention}; hybrids
like \textbf{UNETR} and \textbf{Swin-UNETR} put a transformer encoder on a
convolutional U-Net decoder. \textbf{SAM} (Segment Anything) and its medical
adaptation \textbf{MedSAM} are \emph{foundation models} --- pre-trained on
enormous data and prompted to segment new objects.

\begin{notebox}[A useful dose of scepticism]
A steady stream of papers claims a new Transformer/Mamba architecture beats
U-Net. Recent rigorous benchmarks (``nnU-Net Revisited'', 2024; CNN-vs-Transformer-
vs-Mamba studies, 2025--26) find that when everything is trained and validated
fairly, a well-configured CNN U-Net is still at or near the top, and the fancy
models often cost far more compute for no reliable gain. The takeaway is the same
as Week 3's: \emph{validation discipline beats architectural fashion}. Transformers
shine mainly when you have very large datasets; medical datasets usually are not.
\end{notebox}

% =====================================================================
\section{The medical-imaging twist}
% =====================================================================

Deep learning ``just works'' on natural images partly because ImageNet has
millions of labelled examples. Medical imaging breaks that assumption, and the
differences drive most of the practical decisions:
\begin{itemize}[nosep,leftmargin=*]
  \item \textbf{Small, expensive labels.} Annotation needs a radiologist.
  Augmentation, transfer learning, and self-supervised pre-training are how you
  cope.
  \item \textbf{3D vs 2D.} Many scans are volumes. True 3D nets see context but
  cost memory; 2D-slice nets are cheaper but lose through-plane information.
  \item \textbf{Distribution shift.} A model trained on one scanner/site/vendor
  often degrades on another --- a serious safety issue.
  \item \textbf{Evaluation.} Accuracy is a \emph{terrible} segmentation metric
  (predict ``all background'' and score 99\%). Use \textbf{Dice}, \textbf{IoU},
  and boundary metrics like \textbf{Hausdorff-95}; report on a held-out test set.
\end{itemize}

\begin{notebox}[Worked example: why accuracy lies, and what Dice does]
A tumour occupies $1\%$ of a slice. A model that labels \emph{every} pixel
``background'' is $99\%$ accurate and $100\%$ useless. The \textbf{Dice} score
$\mathrm{Dice}=\dfrac{2|P\cap G|}{|P|+|G|}$ (overlap of prediction $P$ and ground
truth $G$) is $0$ for that model and $1$ only for a perfect mask --- which is why
segmentation is trained and judged with Dice, not accuracy.
\end{notebox}

\begin{notebox}[Region vs boundary: report both]
\textbf{Dice} and \textbf{IoU} $=|P\cap G|/|P\cup G|$ measure region overlap but
say nothing about \emph{boundaries}, so a clinically important edge error can hide
behind a high Dice. \textbf{Hausdorff-95} and average surface distance measure how
far the predicted boundary strays (in mm). Always report a region metric
\emph{and} a boundary metric --- a method can win one and lose the other.
\end{notebox}

\begin{notebox}[Common pitfalls in medical deep learning]
\begin{itemize}[nosep,leftmargin=*]
  \item \textbf{Patient-level leakage.} Slices from the same patient in both train
  and test inflate scores wildly. Split by \emph{patient}, never by slice.
  \item \textbf{Class imbalance.} A tiny lesion drowns in background; use Dice /
  focal losses and report per-class numbers, not a global average.
  \item \textbf{Tuning on the test set.} Pick hyper-parameters on a validation
  split; touch the test set once, at the very end.
  \item \textbf{Distribution shift.} Great numbers on one scanner can collapse on
  another; validate across sites/vendors whenever you can.
\end{itemize}
\end{notebox}

% =====================================================================
\section{The bridge back to the Bayesian picture}
% =====================================================================

A trained CNN is, in our language, a \textbf{learned feature extractor} and a
\textbf{learned prior}. The same network that segments can be read as computing
$p(\text{label}\mid\text{image})$; an encoder--decoder denoiser is doing MAP
estimation with a prior baked into its weights; and in Week 6 a generative
network (VAE/GAN/diffusion) will become an explicit \emph{learned prior} that you
can drop into the reconstruction objective
$\min_x \tfrac12\norm{Gx-y}^2 + \lambda R(x)$ in place of the hand-built $R$. The
plumbing changes; the picture does not.

\begin{keybox}[Do the notebook (and your first project)]
\texttt{notebooks/01\_cnn\_building\_blocks\_from\_scratch.ipynb} (NumPy +
Matplotlib) shows convolution as a feature detector, then trains a tiny network
from scratch with hand-written backprop on synthetic data --- so autodiff,
nonlinearity, and gradient descent are concrete, not magic. For real
\textbf{PyTorch} practice (a CNN classifier, a U-Net), use the cited CS231n /
PyTorch / MONAI tutorials in the \texttt{mic\_soc} environment. This is also where
\textbf{Project 2} (the segmentation project, due in \texttt{week\_6}) starts to
make sense end-to-end.
\end{keybox}

% =====================================================================
\section*{Resources (watch one, read one)}
\addcontentsline{toc}{section}{Resources}
% =====================================================================

\textbf{Foundations (neural nets \& CNNs)}
\begin{itemize}[nosep,leftmargin=*]
  \item 3Blue1Brown, \href{https://www.youtube.com/watch?v=aircAruvnKk}{\emph{But what is a neural network?}} and \href{https://www.youtube.com/watch?v=tIeHLnjs5U8}{backpropagation} (the best visual intuition).
  \item Stanford \href{https://cs231n.github.io/}{CS231n notes} --- ``Convolutional Networks'' and ``Optimisation''.
  \item Dumoulin \& Visin, \href{https://arxiv.org/abs/1603.07285}{\emph{A guide to convolution arithmetic for deep learning}} (stride/padding/dilation, with pictures).
\end{itemize}

\textbf{Hands-on (PyTorch / MONAI)}
\begin{itemize}[nosep,leftmargin=*]
  \item \href{https://pytorch.org/tutorials/beginner/deep_learning_60min_blitz.html}{PyTorch 60-minute blitz}.
  \item \href{https://github.com/Project-MONAI/tutorials}{MONAI tutorials} (medical-imaging PyTorch: 2D/3D segmentation).
\end{itemize}

\textbf{Architectures}
\begin{itemize}[nosep,leftmargin=*]
  \item Ronneberger, Fischer \& Brox (2015), \href{https://arxiv.org/abs/1505.04597}{\emph{U-Net}}; Isensee et al.\ (2021), \href{https://www.nature.com/articles/s41592-020-01008-z}{\emph{nnU-Net}}.
  \item \href{https://arxiv.org/html/2404.09556v2}{\emph{nnU-Net Revisited: A Call for Rigorous Validation}} (2024) --- why configuration beats novelty.
  \item \href{https://link.springer.com/article/10.1007/s00521-026-12198-6}{A comprehensive review of U-Net architectures} (2026); \href{https://www.mdpi.com/2306-5354/12/12/1312}{SAM for medical image segmentation} (review).
\end{itemize}

\textbf{Evaluation \& survey}
\begin{itemize}[nosep,leftmargin=*]
  \item Maier-Hein et al., \href{https://www.nature.com/articles/s41592-023-02151-z}{\emph{Metrics Reloaded}} (choosing the right metric).
  \item Litjens et al.\ (2017), \href{https://www.sciencedirect.com/science/article/pii/S1361841517301135}{\emph{A survey on deep learning in medical image analysis}}.
\end{itemize}

\vfill
\begin{center}\small\itshape
Built for the MIC Summer-of-Code from Prof.\ Suyash P.\ Awate's CS-736
DeepGenerativeModeling deck (CNN half), with attribution. For mentee use on this
project; please do not redistribute the bundled decks outside it.
\end{center}

\end{document}
